\section{Lo que publica cada proveedor}\label{sec:limites:publicado}

Las Tablas~\ref{tab:limites-rates} y~\ref{tab:limites-rates-b} dan, plan a plan, la cuota $Q$ y el \textit{rate} $R$ que declara cada Datasheet, con el semáforo de confianza del Apartado~\ref{sec:propuesta:semaforo}.
Un dato va en verde si es literal del proveedor o lo fijó un panel con cuenta propia o un sondeo, en ámbar si exige una lectura que la fuente no hace literal o si dos fuentes oficiales chocan, y en rojo si el proveedor no lo publica.
Con esos pares las tablas calculan el cociente $R \cdot W / Q$ y el tiempo $t_{\min}$ que tarda el \textit{rate} en agotar la cuota, definidos en la Sección~\ref{sec:limites:geometrica}.
Cada cifra tiene su respaldo en el registro de procedencia de su Datasheet, que guarda para cada valor la URL, la fecha de acceso y la cita literal de la fuente.

Cada proveedor reparte sus términos de consumo, el precio y la información de cómo, cuándo, cuánto y dónde se consume, entre varias páginas: la de precios, la de límites de la API, la guía del relé SMTP o unas preguntas frecuentes.
Los párrafos siguientes recogen los casos en que la fuente de la cifra cambia el cálculo, como el \textit{rate} de Postmark, que solo da la cabecera de la respuesta, o el de MailerSend, que baja con la cuenta en revisión.

\textbf{Una cifra que solo se lee en una cabecera y otra que no se puede observar.}
En Postmark el \textit{rate} de la API REST no está publicado, y la documentación describe el error 429 «Rate Limit Exceeded» sin dar el umbral \citep{postmarkapi}.
La cifra sale de la propia API: los cuatro \textit{endpoints} de envío probados devuelven \texttt{ratelimit-\allowbreak limit:~5000} y \texttt{x-ratelimit-\allowbreak limit-second:~5000}, tanto en la respuesta aceptada (200) como en la rechazada por datos no válidos (422).
Que el límite sea por segundo es una lectura del nombre de la cabecera, no un literal del proveedor.
Un sondeo de 200 peticiones no llegó al 429, y la Datasheet adopta 5.000 peticiones por segundo en ámbar para los cuatro planes, aunque solo se observó en la cuenta gratuita.

SendGrid está en el caso opuesto, con una cifra publicada que no se puede observar.
Da 10.000 peticiones por segundo para \texttt{/v3/mail/send} una sola vez, en las preguntas frecuentes del \textit{endpoint} y sin plan ni ámbito\footnote{\url{https://www.twilio.com/docs/sendgrid/for-developers/sending-email/v3-mail-send-faq}. La cita literal está en el \textit{provenance} de \texttt{03\_datos/apis/sendgrid/datasheet.yaml}.}.
La Datasheet la pone en verde porque es un literal del proveedor.
No hay con qué contrastarla: el envío en modo \textit{sandbox} no devuelve cabeceras \texttt{x-ratelimit-*}.
Ese modo es una opción de la petición que la valida sin entregar el correo, y no tiene que ver con el entorno de pruebas de Amazon SES.

\textbf{Un \textit{rate} que cambia con el estado de la cuenta.}
MailerSend ata su \textit{rate} por minuto a la categoría del \textit{endpoint} y al estado de la cuenta \citep{mailersenddocs}.
Con la reputación sana, \texttt{/v1/email} admite 120 peticiones por minuto en los cuatro planes, es decir, $R = 2$, y \texttt{/v1/bulk-email} entre 10 y 30 según el plan.
Con la cuenta bajo revisión, \texttt{/v1/email} baja a 10 peticiones por minuto y \texttt{/v1/bulk-email} a 1.
El sondeo del Apartado~\ref{subsec:evolucion:p06} midió esos 10 en una cuenta sin aprobar, que es un estado de la suscripción distinto de la revisión por reputación para la que el proveedor publica la cifra.
Las 120 peticiones por minuto están publicadas, pero no observadas.

Cada plan de MailerSend tiene además una cuota diaria de peticiones, común a las dos interfaces de envío, la API REST y el relé SMTP.
Va de 100 peticiones en Free a 500.000 en Professional.
Esa cuota no entra en la Ecuación~\eqref{eq:limites-rates} porque cuenta peticiones y no correos, y la Tabla~\ref{tab:limites-rates} la anota aparte.

\begin{table}[H]
\centering
\scriptsize
\setlength{\tabcolsep}{2.5pt}
\begin{tabular}{|l|r|r|r|c|p{3.2cm}|}
\hline
\textbf{API · plan} & \textbf{$R \cdot W$} & \textbf{$Q$} & \textbf{$R \cdot W / Q$} & \textbf{$t_{\min}$ ($k{=}1$ / $k$ máx)} & \textbf{Datos en ámbar o rojo} \\
\hline
Postmark · Free & $1{,}296 \cdot 10^{10}$ & 100 & 129.600.000 & 0,02 s / 0,0004 s & \textit{rate} medido, no publicado \semambar{} \\
\hline
Postmark · Basic & $1{,}296 \cdot 10^{10}$ & 10.000 & 1.296.000 & 2 s / 0,04 s & \textit{rate} \semambar{}, cuota y ventana \semambar{} \\
\hline
Resend · Free (día) & 864.000 & 100 & 8.640 & 10 s / 0,2 s & ventana diaria \semambar{} \\
\hline
Resend · Free (mes) & 25.920.000 & 3.000 & 8.640 & 300 s / 6 s & ventana mensual \semambar{} \\
\hline
Resend · Pro & 25.920.000 & 50.000 & 518,4 & 83,3 min / 100 s & \textit{rate} por defecto \\
\hline
Resend · Scale & 25.920.000 & 100.000 & 259,2 & 2,78 h / 200 s & \textit{rate} por defecto \\
\hline
Resend · Scale 2,5M & 25.920.000 & 2.500.000 & 10,37 & 69,4 h / 83,3 min & volumen contratado no recogido \\
\hline
Resend · Enterprise & n.c. & Custom & n.c. & n.c. & \textit{rate} \semrojo{} \\
\hline
SendGrid · Trial (día) & 864.000.000 & 100 & 8.640.000 & 0,01 s / 1 petición & \textit{rate} no observable \\
\hline
SendGrid · Essentials 50K & $2{,}592 \cdot 10^{10}$ & 50.000 & 518.400 & 5 s / 0,005 s & \textit{rate} no observable \\
\hline
SendGrid · Pro 100K & $2{,}592 \cdot 10^{10}$ & 100.000 & 259.200 & 10 s / 0,01 s & \textit{rate} no observable \\
\hline
SendGrid · Pro 2,5M & $2{,}592 \cdot 10^{10}$ & 2.500.000 & 10.368 & 250 s / 0,25 s & volumen contratado no recogido \\
\hline
SendGrid · Premier & $2{,}592 \cdot 10^{10}$ & Custom & n.c. & n.c. & cuota \textit{custom}, ventana \semambar{} \\
\hline
MailerSend · Free & 5.184.000 & 500 & 10.368 & 250 s / 4 s & ventana \semambar{}, $k$ \semambar{}, \textit{rate} sin observar. Diaria de 100 pet.: 4 d 50 s ($k{=}1$) \\
\hline
MailerSend · Hobby & 5.184.000 & 5.000 & 1.036,8 & 41,7 min / 36 s & $k$ \semambar{}. Diaria de 1.000 pet.: 4 d 8 min ($k{=}1$) \\
\hline
MailerSend · Starter & 5.184.000 & 50.000 & 103,7 & 6,94 h / 6 min & $k$ \semambar{}, lote de 15/min \semverde{} \\
\hline
MailerSend · Starter 1M & 5.184.000 & 1.000.000 & 5,18 & 5,8 d / 2 h & volumen contratado no recogido. Diaria de 100.000 pet.: 10 d ($k{=}1$) \\
\hline
MailerSend · Professional 2M & 5.184.000 & 2.000.000 & 2,59 & 11,6 d / 4 h & volumen contratado no recogido \\
\hline
MailerSend · Starter 1M, rev. & 432.000 & 1.000.000 & 0,43 & $>$ 30 d / 24 h & \textit{rate} de 10/min medido en la cuenta sin aprobar, volumen contratado no recogido \\
\hline
\end{tabular}
\caption[Relación entre \textit{rate} y cuota en Postmark, Resend, SendGrid y MailerSend]{Relación entre \textit{rate} y cuota en Postmark, Resend, SendGrid y MailerSend. $R \cdot W$ va en peticiones por ventana y $Q$ en correos, con un mes de 30 días. $t_{\min}$ se da con $k = 1$ y con el $k$ máximo de cada API: 50 en Postmark y Resend, 1.000 en SendGrid y 70 en MailerSend, este en ámbar. Cuando la cuota se agota en unas pocas peticiones con el $k$ máximo, $t_{\min}$ se da en número de peticiones en lugar de en segundos. La fila «rev.» usa el \textit{rate} de una cuenta de MailerSend bajo revisión. Un volumen contratado «no recogido» lo publica el proveedor y la Datasheet no lo recoge (Apartado~\ref{subsec:evolucion:p01}). Postmark Pro y Platform tienen la misma cuota y el mismo \textit{rate} que Basic, y se omiten. n.c.: no calculable.}
\label{tab:limites-rates}
\end{table}

\textbf{Categorías de \textit{rate} sin plan asignado.}
Brevo da tres cifras para el mismo \texttt{POST /v3/smtp/email}, un \textit{endpoint} de la API REST pese a su nombre: 1.000 peticiones por segundo en la categoría General, 2.000 en Advanced y 6.000 en Extended \citep{brevodocs}.
La página no cierra qué plan cae en cada una, porque una viñeta reserva Advanced a Professional y Enterprise, y el distintivo «Available in» que encabeza la tabla de Advanced enumera los cinco planes, Free incluido.
Los tres valores van en ámbar, y la Tabla~\ref{tab:limites-rates-b} calcula Professional con la cota inferior, la de General, y Enterprise con la superior, la de Extended.

La guía de envío por lotes da otra cifra para el mismo \textit{endpoint}, «Execute the \textit{endpoint} up to 6000 times per hour», 600 veces menos que los 3.600.000 por hora de la categoría General.
La contradicción sigue abierta, y la tabla calcula con la página de categorías.

\textbf{Cifras que solo valen para el entorno de pruebas.}
En Amazon SES las cifras publicadas son las del entorno de pruebas, que es un estado de la cuenta y no un plan: un correo por segundo, 200 correos cada 24 horas y 50 destinatarios por mensaje \citep{sesdocs}.
Para la cuenta de producción, la página de límites escribe que «this number varies based on your specific use case», y la referencia de la API reserva en \texttt{Max24HourSend} el valor $-1$ para una cuota ilimitada, sin decir qué cuenta lo recibe.
Por eso los cuatro planes que la Datasheet toma de la página de precios \citep{sespricing} reciben la cuota y el \textit{rate} del entorno de pruebas.

Hay dos límites de SES que la Datasheet no pudo modelar con ACM2YAML v0.3.
El primero es el tiempo que SES pide esperar tras un rechazo por limitación, que no tiene campo en el esquema.
La página de errores de la cuota de envío da «up to 10 minutes», y la de problemas del relé SMTP, para las peticiones SMTP rechazadas, una serie de esperas de 5, 10 y 30 segundos seguida de 20 minutos.
La primera página cubre también los rechazos SMTP, así que las dos indicaciones chocan.

El segundo es la petición por segundo que la página de límites aplicaba, cuando se consultó, a todas las operaciones salvo SendEmail, SendRawEmail y SendTemplatedEmail, tres nombres de la versión 1 de la API.
La Datasheet modela la versión 2, y el esquema no podía acotar esa cifra a sus operaciones porque la página no decía cuáles cubría: el envío por lotes no estaba ni entre las exentas ni entre las limitadas.
Un sondeo contra \texttt{GET /v2/email/account} aceptó más de cinco veces ese valor sin fijar el umbral (Apartado~\ref{sec:propuesta:casos}).
Desde finales de octubre de 2026 la página exime SendEmail y SendBulkEmail, y con ese texto la cifra ya se puede acotar.

\textbf{Una cifra publicada y otra medida, con ámbitos distintos.}
Mailgun publica una sola cifra de \textit{rate} para su API REST, «1,500 API requests per second per client IP», en las preguntas frecuentes de su sitio comercial y no en la documentación de desarrollador\footnote{\url{https://www.mailgun.com/ai-info/}. La cita literal está en el \textit{provenance} de \texttt{03\_datos/apis/mailgun/datasheet.yaml}.}.
Leída al pie de la letra es un tope por IP de cliente, un ámbito que el esquema no sabe declarar, y la Datasheet la pone en ámbar como caudal del plan.

Lo medido es otro contador, el de la cabecera del envío aceptado: \texttt{x-ratelimit-limit: 80000}, 80.000 peticiones por ventana.
La ventana la abre la primera petición, y dos mediciones de su duración dieron 9,85 y 9,87 segundos.
La misma respuesta trae \texttt{x-recipient-limit: 2000000}, un contador que limita los destinatarios y no las peticiones.
Es el segundo \textit{rate}, $R_d$, que la Sección~\ref{sec:limites:geometrica} admite junto al de peticiones, y no está en la Datasheet porque el esquema admite un solo \textit{rate} por nodo, sea un plan o un \textit{endpoint}.

La Tabla~\ref{tab:limites-rates-b} calcula con $R = 8.000$ peticiones por segundo, una cifra que no publica Mailgun ni trae la cabecera: es el cociente de esas 80.000 peticiones entre los diez segundos a los que la Datasheet redondea la ventana.
La Datasheet declara las 1.500 peticiones por segundo en el plan y el contador de 80.000 en el \textit{endpoint} de envío.
PRIME4API, que suma los límites del plan a los del \textit{endpoint}, aplica las dos cifras a la vez y calcula con la más baja, las 1.500.

La cuota del plan Free de Mailgun tiene dos cifras: la página de precios da 100 correos al día \citep{mailgunpricing} y el panel de la cuenta muestra 3.000 al mes.
Ninguna de las dos fuentes dice si son dos topes simultáneos, y la tabla calcula con ambos.

\textbf{Un caudal en correos por hora.}
Mailchimp Transactional no publica \textit{rate} en peticiones para el envío, y ni sus páginas ni el sondeo dan cabeceras de \textit{rate} \citep{mailchimpdocs}.
El envío de la cuenta de prueba lo gobierna «a fixed hourly quota of 25 emails», y los correos que pasan de ese tope no se rechazan: entran en una cola que se procesa en la hora siguiente.
La Datasheet interpreta ese tope como un \textit{rate} de 25 correos por hora.
La franquicia de 500 correos de esa cuenta no tiene ventana, y en la cuenta de pago la cuota es un bloque de 25.000 correos \citep{mailchimppricing}, con un \textit{rate} que depende de la reputación y no tiene cifra publicada.

\begin{table}[H]
\centering
\scriptsize
\setlength{\tabcolsep}{2.3pt}
\begin{tabular}{|l|r|r|r|c|p{2.7cm}|}
\hline
\textbf{API · plan} & \textbf{$R \cdot W$} & \textbf{$Q$} & \textbf{$R \cdot W / Q$} & \textbf{$t_{\min}$ ($k{=}1$ / $k$ máx)} & \textbf{Datos en ámbar o rojo} \\
\hline
Brevo · Free (día) & 86.400.000 & 300 & 288.000 & 0,3 s / $<$ 1 pet. & \textit{rate} \semambar{}, hora del reinicio diario \semrojo{} \\
\hline
Brevo · Starter & $2{,}592 \cdot 10^{9}$ & 5.000 & 518.400 & 5 s / 3 pet. & \textit{rate} \semambar{}, $k$ \semambar{}, 99 destinatarios con adjunto \\
\hline
Brevo · Standard & $2{,}592 \cdot 10^{9}$ & 5.000 & 518.400 & 5 s / 3 pet. & ídem Starter, misma cuota y distinto precio \\
\hline
Brevo · Professional & $2{,}592 \cdot 10^{9}$ & 150.000 & 17.280 & 150 s / 75 pet. & categoría Advanced sin atribuir \semambar{} \\
\hline
Brevo · Professional 10M & $2{,}592 \cdot 10^{9}$ & 10.000.000 & 259,2 & 2 h 47 min / 5.000 pet. & volumen contratado no recogido \\
\hline
Brevo · Enterprise & $1{,}5552 \cdot 10^{10}$ & Custom & n.c. & n.c. & categoría Extended \semambar{} \\
\hline
SES · pruebas & 86.400 & 200 & 432 & 200 s / 4 s & ventana de la cuota \semrojo{}, no es cuota de plan \\
\hline
SES · producción & n.c. & no publicada & n.c. & n.c. & cuota \semrojo{} y \textit{rate} \semrojo{} \\
\hline
Mailgun · Free (día) & 691.200.000 & 100 & 6.912.000 & 0,0125 s / $<$ 1 pet. & una o dos cuotas, sin decidir \semambar{}, hora de reposición \semrojo{} \\
\hline
Mailgun · Free (mes) & $2{,}0736 \cdot 10^{10}$ & 3.000 & 6.912.000 & 0,375 s / $<$ 1 pet. & ídem, la diaria por 30 da la mensual \\
\hline
Mailgun · Basic & $2{,}0736 \cdot 10^{10}$ & 10.000 & 2.073.600 & 1,25 s / 0,05 s & ventana del \textit{rate} fuera del esquema \semambar{} \\
\hline
Mailgun · Foundation & $2{,}0736 \cdot 10^{10}$ & 50.000 & 414.720 & 6,25 s / 0,25 s & ídem, \textit{rate} no atribuido por plan \semambar{} \\
\hline
Mailgun · Scale & $2{,}0736 \cdot 10^{10}$ & 100.000 & 207.360 & 12,5 s / 0,5 s & ídem \\
\hline
Mailgun · Scale 2,5M & $2{,}0736 \cdot 10^{10}$ & 2.500.000 & 8.294,4 & 312,5 s / 12,5 s & volumen contratado no recogido \\
\hline
Mailchimp · prueba & en correos/h & 500 & no formable & 20 h (cualquier $k$) & cuota sin ventana \semrojo{}, $k$ sin tope \semrojo{} \\
\hline
Mailchimp · bloque de pago & sin cifra & 25.000 & no formable & n.c. & \textit{rate} sin publicar, ventana del bloque \semambar{} \\
\hline
\end{tabular}
\caption[Relación entre \textit{rate} y cuota en Brevo, Amazon SES, Mailgun y Mailchimp]{Relación entre \textit{rate} y cuota en Brevo, Amazon SES, Mailgun y Mailchimp, con las convenciones de la Tabla~\ref{tab:limites-rates}. El $k$ máximo vale 2.000 en Brevo (en ámbar), 50 en SES y 1.000 en Mailgun. El \textit{rate} de Mailgun es el de la cabecera del envío, 80.000 peticiones en una ventana medida de unos diez segundos. Su $R = 8.000$ es el cociente de esas 80.000 peticiones entre los diez segundos a los que la Datasheet redondea la ventana, y no una cifra publicada. Su $t_{\min}$ con $k$ máximo lo gobierna el contador de 2.000.000 de destinatarios sobre la misma ventana. Las dos filas de SES son estados de la cuenta y valen para los cuatro planes. Con adjunto, Brevo baja el tope a 99 destinatarios por petición. Mailchimp declara su \textit{rate} en correos por hora, y por eso no tiene $R \cdot W$. n.c.: no calculable.}
\label{tab:limites-rates-b}
\end{table}

\textbf{La especificación formal dice menos que la prosa.}
La especificación formal de una API es su descripción legible por máquina: un fichero OpenAPI o Swagger que enumera cada ruta con sus parámetros y sus respuestas.
La de SendGrid, en el repositorio \texttt{twilio/sendgrid-oai}, no menciona los \textit{rates} ni el tope de \texttt{to}, y el tope de 1.000 destinatarios por petición, que suma \texttt{to}, \texttt{cc} y \texttt{bcc}, solo aparece en la referencia en prosa \citep{sendgriddocs}.
La de MailerSend, publicada en SwaggerHub, tampoco declara el tope de \texttt{to} ni la respuesta 429 \citep{mailersenddocs}.
La de Brevo, una definición Swagger de 207 rutas, no menciona el \textit{rate} ni una vez \citep{brevodocs}.

Amazon SES publica sus límites en cuatro documentos: la página de límites, la referencia de la operación de envío, la especificación formal de la versión 2 y la guía del relé SMTP \citep{sesdocs}.
En una tabla de doce límites por esos cuatro documentos, de la cuota diaria a la espera tras el rechazo, quedan vacías 32 de las 48 celdas, y la columna de la especificación formal lo está entera.
Una Datasheet construida solo desde la especificación formal recogería menos límites que la construida desde la prosa en esas cuatro APIs, y la comparación no se hizo en las otras cuatro.
